\section{Method}

\subsection{Tree-KV state and copy-on-write forks}

For each transformer layer, the KV pool stores fixed-size pages with shape
$[P,S,H_{kv},D]$, where $P$ is the number of physical pages and the default page
size $S$ is 16 tokens.  A branch records its parent, its root-to-node token
count, and an ordered block table of physical page identifiers.  Forking a
branch increments page references and copies no covered page.  When a child
writes to a shared page, the pool first copies that page and redirects only the
writer.  Thus branch creation is cheap while divergence pays the copy cost
lazily.

Pruning is defined only for leaves.  It decrements every covered page reference
and immediately returns zero-reference pages to the allocator.  Full frozen
pages may additionally enter a content-addressed deduplication pass: pages with
identical raw bytes are redirected to one physical page and duplicates are
freed.  The pool exports logical tokens, physical tokens, used pages, and the
reuse multiplier
\begin{equation}
  R_{KV}=\frac{\text{logical token references}}
                 {\text{physical tokens stored}}.
\end{equation}
The served wire contract requires a positive physical count and $R_{KV}\geq1$.

Convergent branch merging is a scheduler-level operation distinct from the
deduplication pass.  The wire protocol defines a terminal branch-merge event and
a live target branch.  The present engine does not yet drive those merge events;
it implements pool-level full-page deduplication.  This separation prevents the
paper from treating a storage primitive as an already complete semantic merge
algorithm.

\subsection{Branch-aware tree attention}

At a decode step, each active branch contributes one query.  Its block table
and context length select exactly the root-to-node K/V path.  A query cannot
attend to tokens owned only by a sibling, which is the tree analogue of a causal
mask.  The normative reference implementation uses PyTorch and fp32
accumulation.  An import-guarded Triton decode path is intended to match that
reference within dtype-specific tolerances; fused prefill remains future work.

\subsection{Rust scheduling policies}

The scheduler owns a branch tree whose root begins active.  Engine events report
sampled tokens, branch exhaustion, and optional external value scores.  Commands
authorize exactly one continuation step, fork an active branch, kill it with a
stable reason, or finalize it.  Beam search retains the highest-ranked live
frontier, best-first selects the strongest available branch, and MCTS balances
estimated value with exploration.  Deterministic random-number handling makes
policy traces replayable for a fixed configuration and seed.

External scoring is synchronized with decoding: once a branch requests a value,
it cannot consume another token until the corresponding score arrives.  This
prevents a fast branch from outrunning the policy decision that is supposed to
govern it.

\subsection{Budget controller}

Every sampled token consumes one unit of the total tree budget.  Per-branch and
global exhaustion produce explicit kill reasons, and the engine reclaims killed
leaf state immediately.  The budget is therefore a hard execution invariant,
not a post-hoc accounting field.  A future thinking-effort adapter can map a
model-level control to width, depth, scorer cadence, and token budget while
retaining this invariant.
